\documentclass[11pt,a4paper]{article}
\usepackage[utf8]{inputenc}
\usepackage[margin=1in]{geometry}
\usepackage{amsmath, amssymb, amsfonts}
\usepackage{booktabs}
\usepackage{hyperref}
\usepackage{titlesec}
\usepackage{microtype}

% Title Configuration
\title{\textbf{Comparative Analysis of Generative Frameworks:\\ Focus on DDPM vs. Flow Matching vs. Rectified Flow}}
\author{\textbf{Presenters:} Joachim Hofmann \& Alex Hanselmann \\
\small \textbf{Context:} Master's Course in Generative Machine Learning}
\date{\textbf{Date:} May 28, 2026}

\begin{document}

\maketitle

\hrule
\vspace{0.4cm}

\section{Introduction and Research Objectives}
This report documents the findings of our research within Bundle A2 of the Generative Machine Learning curriculum[cite: 87]. We address the recent paradigm shift in generative modeling: the move from stochastic diffusion processes to deterministic flow-based methods[cite: 88]. While Denoising Diffusion Probabilistic Models (DDPM) established the baseline for high-fidelity synthesis, their reliance on stochastic reverse chains introduces significant computational bottlenecks[cite: 89].

Our primary mission is a comparative evaluation of DDPM, Flow Matching (FM), and Rectified Flow (RF) using the CelebA-64 dataset[cite: 90]. We aim to rigorously analyze how these frameworks handle the transport from a Gaussian source to a data distribution, focusing on two research goals[cite: 91]:
\begin{itemize}
    \item \textbf{Minimization of Operational Cost:} We define this through ``NFE'' (Number of Function Evaluations)---the count of neural network passes required for a single sample---and ``Wall-Clock Cost'' (Latency), the physical execution time in milliseconds (ms)[cite: 91].
    \item \textbf{Mathematical Formulation Stability:} This refers to the geometric precision of the generative trajectories[cite: 92]. We quantify stability using the Mean Squared Error (MSE), which measures the ``Path Deviation'' of the model's predicted trajectory from a theoretical straight line[cite: 93].
\end{itemize}
By analyzing these metrics, we seek to demonstrate how path ``straightening'' directly impacts hardware efficiency and model stability[cite: 94].

\section{Theoretical Framework: Stochastic vs. Deterministic Transport}
The architectural divergence between these frameworks is rooted in the choice between Stochastic Differential Equations (SDE) and Ordinary Differential Equations (ODE)[cite: 95]. DDPM models the generative process as a reverse-time SDE, incorporating ``stochastic jitter'' to remain stable[cite: 96]. Conversely, Flow Matching and Rectified Flow utilize an ODE framework to define a smooth, deterministic velocity field[cite: 97].

\begin{table}[h]
\centering
\caption{Systematic Comparison of Generative Paradigms}
\label{tab:paradigms}
\vspace{0.2cm}
\small
\begin{tabular}{lll}
\toprule
\textbf{Dimension} & \textbf{Diffusion Models (DDPM)} & \textbf{Rectified Flow (RF)} \\
\midrule
Systemic Basis & SDE: $dx_t = f(t)x_t dt + g(t)dw_t$ & ODE: $dx_t = v(x_t, t)dt$ \\
\addlinespace
Path Interpolation & Non-linear: $X_t = \sqrt{\bar{\alpha}_t} X_1 + \sqrt{1 - \bar{\alpha}_t} \epsilon$ & Linear: $X_t = (1 - t)X_0 + tX_1$ \\
\addlinespace
Learning Target & Added Noise ($\epsilon$) & Velocity Field ($v$) \\
\addlinespace
Loss Function & $\mathcal{L} = \mathbb{E} [ \| \epsilon - \hat{\epsilon}_\theta \|^2 ]$ & $\mathcal{L} = \mathbb{E} [ \| (X_1 - X_0) - v_\theta \|^2 ]$ \\
\addlinespace
Inference Logic & Stochastic Markov Chain & Deterministic Euler Step \\
\bottomrule
\end{tabular}
\end{table}

\subsection{Path Entanglement and the Mean Value Problem}
A fundamental challenge in high-dimensional transport is ``Path Entanglement.'' Even with a linear interpolation between $X_0$ and $X_1$, different pairs of noise and data often result in crossing trajectories[cite: 100]. Since a neural network $v_\theta(x, t)$ is a single-valued function, it is mathematically forced to learn the average (mean) of all target directions at any intersection point[cite: 101]. This forcing leads to a time-variant, curved velocity field, which manifests empirically as ``Blur'' in generated samples and high inference costs due to the need for many small integration steps[cite: 102].

\subsection{The Mathematical Link: Eulerian vs. Lagrangian}
Flow Matching and Rectified Flow describe the same physical transport but through different lenses[cite: 103]. Flow Matching typically adopts an Eulerian (Field-centric) perspective, focusing on the continuity equation[cite: 104]:
\begin{equation}
\frac{\partial \rho_t(x)}{\partial t} + \nabla \cdot (\rho_t(x)v(x, t)) = 0
\end{equation}
to preserve probability mass[cite: 104]. Rectified Flow adopts a Lagrangian (Particle-centric) perspective, focusing on the trajectories of individual random variable paths $X_t$[cite: 105]. The goal of Rectified Flow is to solve the entanglement problem by explicitly straightening these paths[cite: 106].

\section{The Rectified Flow Framework: From Chaos to Straight Highways}
Rectified Flow provides a structured, three-phase approach to eliminate path curvature, enabling extreme inference efficiency[cite: 107].

\begin{enumerate}
    \item \textbf{Phase 1: Initial Training/Rectification (1-RF):} The model learns a velocity field $v_\theta$ to minimize the least-squares error along linear paths[cite: 108]. While this establishes a valid flow, the paths still cross and curve because $X_0$ and $X_1$ are coupled independently[cite: 109].
    \item \textbf{Phase 2: The Reflow Step (2-RF):} This is the critical ``Flow-Rewiring'' phase[cite: 110]. By simulating the ODE from Phase 1, we find the deterministic endpoints $X^{\text{rect}}_1$ for each $X_0$[cite: 111]. This process re-pairs the noise and data to minimize the quadratic transport cost[cite: 112]:
    \begin{equation}
    \min_\pi \int \|x_0 - x_1\|^2 d\pi(x_0, x_1)
    \end{equation}
    The result is a set of non-intersecting, perfectly straight trajectories[cite: 112].
    \item \textbf{Phase 3: Jump:} Once the paths are straightened and the velocity field becomes time-independent ($v(x, t) \approx v(x, 0)$), a student model $\alpha_\phi(X_0)$ is trained to predict the exact ODE endpoint $X_1$ in a single step, effectively skipping the simulation path entirely[cite: 113].
\end{enumerate}

\noindent\textbf{Analysis:} The ``Reflow'' step is what allows the generative process to collapse from many steps to one[cite: 114]. By minimizing the quadratic transport cost, the model regularizes the velocity field toward homogeneity[cite: 115]. This removes the necessity of an iterative solver to compensate for direction changes during inference[cite: 116].

\section{Implementation Architecture and Evaluation Metrics}
To ensure a rigorous ``fair-comparison'' environment, we utilized a routing node architecture where all models originate from a shared initial noise distribution $\mathcal{N}(0, \mathbf{I})$[cite: 117].

\subsection{Evaluation Architecture}
\begin{itemize}
    \item \textbf{Routing Node:} Distributes the shared latent noise to three execution paths: the Stochastic Reverse Chain (DDPM), the Deterministic ODE Solver (FM), and the 1-Step Jump Network (RF)[cite: 118].
    \item \textbf{Backbone Integration:} We employ a convolutional U-Net conditioned on the time variable $t$[cite: 119]. The scalar $t$ is first transformed via sinusoidal position encoding[cite: 120].
    \item \textbf{MLP Specifics:} The embedding is processed by a Multi-Layer Perceptron using a specific \texttt{nn.Sequential(nn.Linear(...), nn.ReLU())} architecture[cite: 121]. The output is projected and added to the U-Net feature maps, allowing the network to dynamically adjust to the current timestep[cite: 122].
\end{itemize}

\subsection{Measurement Layer}
We evaluate the performance across three specific metrics[cite: 123]:
\begin{itemize}
    \item \textbf{NFE $\downarrow$:} The total number of network executions per image[cite: 123].
    \item \textbf{Latency (ms) $\downarrow$:} Physical wall-clock time on hardware[cite: 124].
    \item \textbf{MSE (Path Deviation) $\downarrow$:} The mean squared error measuring drift from a linear trajectory[cite: 124].
\end{itemize}

\end{document}